% from ega2-1-fr.tex
\oldpage[II]{127}

éléments inversibles de $\Gamma(U,\mathcal{A})$ (resp.
$\Gamma(U,\mathcal{B})$) pour tout ouvert $U\subset X$~; ce sont des faisceaux
de groupes multiplicatifs, et $N_{\mathcal{B}/\mathcal{A}}$, restreint à
$\mathcal{B}^*$, est un \emph{homomorphisme}
$\mathcal{B}^*\to\mathcal{A}^*$ de faisceaux de groupes
(\hyperref[II.6.5.1-fr]{6.5.1}). Soit $\mathfrak{L}$ l'ensemble des couples
$(U_\lambda,\eta_\lambda)$, ayant la propriété suivante~: $U_\lambda$ est un
ouvert de $X$, et $\eta_\lambda$ est un isomorphisme
$\mathcal{L}'|U_\lambda\simeq\mathcal{B}|U_\lambda$ de
$(\mathcal{B}|U_\lambda)$-Modules. Par hypothèse, les $U_\lambda$ forment un
recouvrement de $X$~; pour deux indices quelconques $\lambda$, $\mu$, on pose
$\omega_{\lambda\mu}=(\eta_\lambda|U_\lambda\cap U_\mu)
\circ(\eta_\mu|U_\lambda\cap U_\mu)^{-1}$, automorphisme de
$\mathcal{B}|U_\lambda\cap U_\mu$, qui s'identifie canoniquement à une section
de $\mathcal{B}^*$ au-dessus de $U_\lambda\cap U_\mu$, et
$(\omega_{\lambda\mu})$ est un $1$-cocycle du recouvrement
$\mathfrak{U}=(U_\lambda)$ à valeurs dans $\mathcal{B}^*$
(\hyperref[0.5.4.7-fr]{0, 5.4.7}). Le fait que
$N_{\mathcal{B}/\mathcal{A}}:\mathcal{B}^*\to\mathcal{A}^*$ est un
homomorphisme entraîne alors que
$(N_{\mathcal{B}/\mathcal{A}}\omega_{\lambda\mu})$ est un $1$-cocycle de
$\mathfrak{U}$ à valeurs dans $\mathcal{A}^*$, et il lui correspond donc (à un
isomorphisme unique près) un $\mathcal{A}$-Module inversible $\mathcal{L}$
(\hyperref[0.5.4.7-fr]{0, 5.4.7}) que nous désignerons par
$N_{\mathcal{B}/\mathcal{A}}(\mathcal{L}')$ et appellerons \emph{la norme du
$\mathcal{B}$-Module inversible $\mathcal{L}'$}.

Soit $\mathfrak{M}$ une partie de $\mathfrak{L}$ telle que les $U_\lambda$
correspondants forment encore un recouvrement de $X$, et soit $\mathfrak{B}$
ce recouvrement~; la restriction du cocycle $(\omega_{\lambda\mu})$ à
$\mathfrak{B}$ définit un $1$-cocycle
$(N_{\mathcal{B}/\mathcal{A}}\omega_{\lambda\mu})$ de $\mathfrak{B}$ à
valeurs dans $\mathcal{A}^*$, restriction du $1$-cocycle
$(N_{\mathcal{B}/\mathcal{A}}\omega_{\lambda\mu})$ de $\mathfrak{U}$~; il est
clair qu'il y a un isomorphisme canonique du $\mathcal{A}$-Module inversible
défini par ce $1$-cocycle de $\mathfrak{B}$ sur
$N_{\mathcal{B}/\mathcal{A}}(\mathcal{L}')$, ce qui permet de définir ce
$\mathcal{A}$-Module inversible par un sous-recouvrement quelconque de
$\mathfrak{U}$. Cette possibilité montre aussitôt que, si
$\mathcal{L}'_1$, $\mathcal{L}'_2$ sont deux $\mathcal{B}$-Modules inversibles,
on a, en vertu de (\hyperref[II.6.5.1.1-fr]{6.5.1.1}) et
(\hyperref[II.6.5.1.2-fr]{6.5.1.2}),
\[
\label{II.6.5.2.1-fr}
N(\mathcal{L}'_1\otimes_{\mathcal{B}}\mathcal{L}'_2)
=N(\mathcal{L}'_1)\otimes_{\mathcal{A}}N(\mathcal{L}'_2)
\tag{6.5.2.1}
\]
et
\[
\label{II.6.5.2.2-fr}
N_{\mathcal{B}/\mathcal{A}}(\mathcal{B})=\mathcal{A}
\tag{6.5.2.2}
\]
ainsi que
\[
\label{II.6.5.2.3-fr}
N((\mathcal{L}')^{-1})=(N(\mathcal{L}'))^{-1}
\tag{6.5.2.3}
\]
à des isomorphismes canoniques près. En outre il résulte de
(\hyperref[II.6.5.1.3-fr]{6.5.1.3}) que si $\mathcal{L}$ est un
$\mathcal{A}$-Module inversible et si $\mathcal{B}$ est de rang constant $n$
sur $\mathcal{A}$ dans le cas (I) (resp.
$\mathcal{B}\otimes_{\mathcal{A}}\mathcal{R}(X)$ de rang constant $n$ sur
$\mathcal{R}(X)$ dans le cas (II)), on a, à un isomorphisme canonique près
\[
\label{II.6.5.2.4-fr}
N_{\mathcal{B}/\mathcal{A}}(\mathcal{L}\otimes_{\mathcal{A}}\mathcal{B})
=\mathcal{L}^{\otimes n}.
\tag{6.5.2.4}
\]
\end{env}

\begin{env}[6.5.3]
\label{II.6.5.3-fr}
Montrons maintenant que $N_{\mathcal{B}/\mathcal{A}}(\mathcal{L}')$ est un
\emph{foncteur covariant} dans la catégorie des $\mathcal{B}$-Modules
inversibles. Soit en effet $h':\mathcal{L}'_1\to\mathcal{L}'_2$ un
homomorphisme de $\mathcal{B}$-Modules inversibles, et soit
$\mathfrak{B}=(U_\lambda)$ un recouvrement ouvert de $X$ tel que pour tout
$\lambda$, on ait des $(\mathcal{B}|U_\lambda)$-isomorphismes
$\eta^{(1)}_\lambda:\mathcal{L}'_1|U_\lambda\simeq\mathcal{B}|U_\lambda$ et
$\eta^{(2)}_\lambda:\mathcal{L}'_2|U_\lambda\simeq\mathcal{B}|U_\lambda$~;
il y a donc pour chaque $\lambda$ un endomorphisme $h'_\lambda$ de
$\mathcal{B}|U_\lambda$ tel que
$h'_\lambda\circ\eta^{(1)}_\lambda
=\eta^{(2)}_\lambda\circ(h'|U_\lambda)$, et on peut évidemment identifier
$h'_\lambda$ à une section de $\mathcal{B}$ au-dessus de $U_\lambda$
(\hyperref[0.5.1.1-fr]{0, 5.1.1}). Donc, pour tout couple $(\lambda,\mu)$
d'indices, les restrictions à $U_\lambda\cap U_\mu$ de
$(\eta^{(2)}_\lambda)^{-1}\circ h'_\lambda\circ\eta^{(1)}_\lambda$ et de
$(\eta^{(2)}_\mu)^{-1}\circ h'_\mu\circ\eta^{(1)}_\mu$ coïncident. On en
déduit pour les $1$-cocycles $(\omega^{(1)}_{\lambda\mu})$ et
$(\omega^{(2)}_{\lambda\mu})$ à valeurs dans $\mathcal{B}^*$ correspondant à
$\mathcal{L}'_1$ et $\mathcal{L}'_2$, la relation
\[
\omega^{(2)}_{\lambda\mu}h'_\mu
=h'_\lambda\omega^{(1)}_{\lambda\mu}.
\]

